\section{지수의 확장}

2026-08-23.

\subsection{유리수 지수}

\(n^{1/k}\)는 \(k\)번 곱해서 \(n\)이 되는 양수, 즉 \(\sqrt[k]{n}\)이다.

\[
n^{1/k}=\sqrt[k]{n},\qquad n^{m/k}=\big(\sqrt[k]{n}\big)^m=\sqrt[k]{n^m}\qquad(n>0)
\]

예: \(8^{1/3}=2\), \(8^{2/3}=4\).

\subsection{무리수 지수}

\(n^{\sqrt k}\)에서 \(\sqrt k\)는 (완전제곱이 아니면) 무리수라서 ``몇 번
곱한다''로는 정의할 수 없다. 두 가지 방법으로 정의한다.

\begin{enumerate}
\def\labelenumi{\arabic{enumi}.}
\tightlist
\item
  유리수 근삿값의 극한: \(\sqrt2=1.4,\ 1.41,\ 1.414,\dots\)이면
  \(2^{1.4},2^{1.41},2^{1.414},\dots\)가 한 값으로 수렴한다.
\item
  지수·로그 함수로:
\end{enumerate}

\[
n^{x}=e^{x\ln n}\qquad\Rightarrow\qquad n^{\sqrt k}=e^{\sqrt k\,\ln n}
\]

예: \(2^{\sqrt2}=e^{\sqrt2\ln2}\approx2.6651\).

참고: \(2^{\sqrt2}\)는 겔폰트--슈나이더 정리에 의해 초월수다.
